\subsection{可分解包络}

\(\mathfrak{gl}(V)\) 的一族可分解李子代数的交显然是可分解的。因而，若 g 是 \(\mathfrak{gl}(V)\) 的一个李子代数，则 \(\mathfrak{gl}(V)\) 的包含 g 的可分解李子代数所成的集有一个最小元素，称为 g 的可分解包络；在本节中，这个包络将记作 \(e(\mathfrak{g})\)。

命题 4. 设 \(\mathfrak{g}\) 是 \(\mathfrak{gl}(\mathrm{V})\) 的一个李子代数，\(\mathfrak{n}\) 是 \(\mathfrak{g}\) 的一个理想。则 \(\mathfrak{n}\) 与 \(e(\mathfrak{n})\) 都是 \(e(\mathfrak{g})\) 的理想，且 \([e(\mathfrak{g}), e(\mathfrak{n})] = [\mathfrak{g}, \mathfrak{n}]\)。

设 \(\mathfrak{g}_1\) 是那些 \(x\in \mathfrak{gl}(\mathrm{V})\)（满足 \([x,\mathfrak{n}]\subset [\mathfrak{g},\mathfrak{n}]\)）所成的集。这是 \(\mathfrak{gl}(\mathrm{V})\) 的一个可分解李子代数，包含 \(\mathfrak{g}\)，从而包含 \(e(\mathfrak{g})\)，参见第 1 节命题 3；换言之，\([e(\mathfrak{g}),\mathfrak{n}]\subset [\mathfrak{g},\mathfrak{n}]\)。设 \(\mathfrak{n}_1\) 是那些 \(y\in \mathfrak{gl}(\mathrm{V})\) 所成的集，使得

\[
[ e (\mathfrak {g}), y ] \subset [ \mathfrak {g}, \mathfrak {n} ].
\]

由前所述，这是 \(\mathfrak{gl}(V)\) 的一个包含 n 的可分解李子代数，从而包含 \(e(\mathfrak{n})\)；换言之，\([e(\mathfrak{g}), e(\mathfrak{n})] \subset [\mathfrak{g}, \mathfrak{n}]\)，于是

\[
[ e (\mathfrak {g}), e (\mathfrak {n}) ] = [ \mathfrak {g}, \mathfrak {n} ].
\]

由此推出 \([e(\mathfrak{g}),\mathfrak{n}]\subset [e(\mathfrak{g}),e(\mathfrak{n})]\subset \mathfrak{n}\)，于是 \(\mathfrak{n}\) 与 \(e(\mathfrak{n})\) 都是 \(e(\mathfrak{g})\) 的理想。

推论 1. (i) \(\mathcal{D}^i\mathfrak{g} = \mathcal{D}^i e(\mathfrak{g})\) 对 \(i\geq 1\) 成立，且 \(\mathcal{C}^i\mathfrak{g} = \mathcal{C}^i e(\mathfrak{g})\) 对 \(i\geq 2\) 成立。

\begin{enumerate}
\def\labelenumi{(\roman{enumi})}
\setcounter{enumi}{1}
\tightlist
\item
  若 g 交换（相应地，幂零、可解），则 \(e(\mathfrak{g})\) 交换（相应地，幂零、可解）。
\end{enumerate}

断言 (i) 由命题 4 对 i 用归纳法得到，而 (ii) 由 (i) 得到。

推论 2. 设 \(\mathfrak{r}\) 是 \(\mathfrak{g}\) 的根基。若 \(\mathfrak{g}\) 可分解，则 \(\mathfrak{r}\) 可分解。

事实上，由命题 4 与推论 1，\(e(\mathfrak{r})\) 是 \(\mathfrak{g}\) 的一个可解理想，故 \(e(\mathfrak{r}) = \mathfrak{r}\)。

